\begin{frame}
	\frametitle{Ответы на замечания Обронова~И.В.}
	\textbf{Замечание}: Во второй главе предложен способ подавления ММЧВС с помощью модового фильтра. Показано, что применение фильтра позволяет увеличить допустимую длину волокна доставки с 2.5 до 3.0 м. Однако автором отмечено, что на длине 4 м подавление ММЧВС приводит к развитию ВКР, что, фактически, компенсирует положительный эффект. Каким образом предполагается решать проблему конкуренции нелинейных процессов при дальнейшем масштабировании длины волокна или мощности, или же достигнутая длина является принципиальным пределом для данного типа волокон?
	
	\textbf{Ответ:} Достигнутая длина 3~м~--- это техническое ограничение использованного способа подавления эффекта (модового фильтра) и применённого волокна доставки. Для дальнейшего масштабирования оптической мощности представляется возможным рассмотреть следующие варианты:
	\begin{enumerate}
		\item Использование LMA-волокна доставки с адиабатическим увеличением диаметра сердцевины.
		\item Использование фосфоросиликатного ОВ (большая величина стоксова сдвига, вне области усиления Yb).
		\item Спектральная фильтрация ВКР с помощью FBG-решёток.
	\end{enumerate}
\end{frame}


\begin{frame}
	\frametitle{Ответы на замечания Обронова~И.В.}
	\textbf{Замечание}: Четвертая глава посвящена повышению точности метода пьезоэлектрической резонансной лазерной калориметрии. В качестве объекта исследования выбран кристалл LBO, однако во введении и первой главе в качестве материала для генерации излучения среднего ИК-диапазона фигурирует кристалл PPLN. Насколько разработанные поправочные выражения и предложенная методика коррекции результатов применимы к данному кристаллу?
	
	\textbf{Ответ}: Поправочные выражения универсальны и зависят только от теплофизических параметров образца и условий эксперимента, а не от материала. Поэтому методика применима и к PPLN. LBO выбран как модельный объект, так как у него низкое поглощение и большие размеры, из-за чего погрешности от нестабильности температуры и теплоотдачи проявляются максимально сильно.
\end{frame}